\documentclass[10pt]{article}
\usepackage[T1]{fontenc}
\usepackage[utf8]{inputenc}
\usepackage[french]{babel}
\usepackage{amsmath,amssymb,mathrsfs,mathpazo}
\usepackage{geometry}
\geometry{a4paper,margin=2cm}
\usepackage{graphicx,array,tabularx,tabularray,colortbl,multirow,multicol,booktabs,enumitem,pifont,ulem,fancybox,diagbox,cancel}
\usepackage{pstricks,pstricks-add,pst-plot,pst-tree,pst-node,pst-text,pst-func,pst-math,pst-bspline,pst-3dplot,pst-eucl,pst-solides3d}
\let\pscalculate\undefined
\usepackage{pst-bezier}
\usepackage{tikz,pgfplots,tkz-tab}
\usetikzlibrary{arrows,arrows.meta,patterns,calc,positioning,trees,decorations.markings,shapes.geometric,shapes.misc,shapes.arrows,intersections,decorations.pathmorphing,decorations.pathreplacing,backgrounds,shadows,babel}
\pgfplotsset{compat=1.18}
\usepgfplotslibrary{fillbetween}
\usepackage[np]{numprint}
\usepackage{listings,makecell,pas-tableur,twoopt}
\let\barre\undefined
\usepackage[pstricks]{tablvar}
\definecolor{darkgreen}{RGB}{0,100,0}
\providecommand{\red}{\color{red}}
\providecommand{\blue}{\color{blue}}
\providecommand{\green}{\color{green}}
\providecommand{\black}{\color{black}}
\providecommand{\white}{\color{white}}
\providecommand{\gray}{\color{gray}}
\providecommand{\cyan}{\color{cyan}}
\providecommand{\magenta}{\color{magenta}}
\providecommand{\yellow}{\color{yellow}}
\providecommand{\violet}{\color{violet}}
\providecommand{\orange}{\color{orange}}
\providecommand{\euro}{\texteuro}
\providecommand{\EUR}{\texteuro}
\providecommand{\ieme}{\textsuperscript{e}}
\providecommand{\degres}{\ensuremath{^\circ}}
\providecommand{\degree}{\ensuremath{^\circ}}
\providecommand{\vv}[1]{\overrightarrow{#1}}
\providecommand{\Coord}[1]{#1}
\providecommand{\decosix}{}
\providecommand{\starredbullet}{\ensuremath{\bullet}}
\pagestyle{empty}
\setlength{\parindent}{0pt}
\usepackage{adjustbox}\DeclareUnicodeCharacter{2236}{:}\DeclareUnicodeCharacter{2190}{\ensuremath{\leftarrow}}
\providecommand{\R}{\mathbb{R}}
\providecommand{\N}{\mathbb{N}}
\providecommand{\Z}{\mathbb{Z}}
\providecommand{\Q}{\mathbb{Q}}
\providecommand{\D}{\mathbb{D}}
\providecommand{\C}{\mathbb{C}}
\providecommand{\e}{{\mathrm{e}}}
\providecommand{\vect}[1]{\overrightarrow{#1}}
\providecommand{\vectt}[1]{\overrightarrow{\mathrm{#1}}}
\providecommand{\vvt}[1]{\overrightarrow{\mathrm{#1}}}
\providecommand{\np}[2][]{{#2~#1}}
\providecommand{\vv}[1]{{\overrightarrow{#1}}}
\renewcommand{\d}{{\mathrm{d}}}
\providecommand{\degres}{{^\circ}}
\providecommand{\degre}{{^\circ}}
\providecommand{\blue}{\color{blue}}
\providecommand{\red}{\color{red}}
\providecommand{\enskip}{\quad}
\providecommand{\llbracket}{\lbrack\!\lbrack}
\providecommand{\rrbracket}{\rbrack\!\rbrack}
\providecommand{\euro}{{\text{€}}}
\providecommand{\no}{{\text{n°}}}
\providecommand{\interval}[2]{{[#1;#2]}}
\providecommand{\intervFO}[2]{{[#1;#2[}}
\providecommand{\dotfill}{\ldots}

\providecommand{\ds}{\displaystyle}
\providecommand{\textding}[1]{#1}
\providecommand{\Oij}{\(\mathrm{O};\vec{\imath},\vec{\jmath}\)}
\providecommand{\Oijk}{\(\mathrm{O};\vec{\imath},\vec{\jmath},\vec{k}\)}
\providecommand{\Ouv}{\(\mathrm{O};\vec{u},\vec{v}\)}
\providecommand{\intervFF}[2]{\ensuremath{[#1;#2]}}
\renewcommand{\euro}{\ifmmode\text{\texteuro}\else\texteuro\fi}
\DeclareUnicodeCharacter{00B0}{\ifmmode^{\circ}\else\textdegree\fi}
\DeclareUnicodeCharacter{20AC}{\ifmmode\text{\texteuro}\else\texteuro\fi}
% Énoncés : Éducation nationale. Transcription et corrigés : APMEP, Denis Vergès et contributeurs. Découpage et classement : CoopMaths, Sébastien Lozano, Éric Elter et contributeurs. Adaptation de la présentation : Maths Explorer.
% CC BY-SA : https://coopmaths.fr/www/about/
% https://coopmaths.fr/alea/static/eam/2026/tex/eam_technologique_2026_06_metropole_2.tex
\begin{document}
\noindent\textit{APMEP / CoopMaths. CC BY-SA. Présentation adaptée par Maths Explorer.}\par\medskip
Une salle de sport propose deux offres d'abonnement :

Abonnement n°1 : 250 \euro{} en 2026 puis une augmentation de 30 \euro{} par an.

Abonnement n°2 : 200 \euro{} en 2026 puis une augmentation de 10\% par an.

\textbf{Les parties A et B sont indépendantes.}


\subsection*{PARTIE A : Abonnement n°1.}

Pour tout entier naturel $n$, on note $a_n$ le montant de l'abonnement n°1 pour l'année $2026+n$.

On a ainsi $a_0=250$.
\begin{enumerate}
\item
\begin{enumerate}[label=\alph*)]
\item Calculer le montant de l'abonnement pour 2027.
\item Calculer $a_2$ et interpréter le résultat dans le contexte de l'exercice.
\end{enumerate}
\item Exprimer, pour tout entier naturel $n$, $a_{n+1}$ en fonction de $a_n$.
\item En déduire la nature de la suite $(a_n)$ et 
préciser sa raison.
\end{enumerate}

\bigskip

\subsection*{PARTIE B : Abonnement n°2.}

Pour tout entier naturel $n$, on note $b_n$ le montant de l'abonnement n°2 pour l'année $2026+n$.

On a ainsi $b_0=200$.
\begin{enumerate}
\item Justifier que le montant de l'abonnement en 2027 s'élève à 220 \euro{}.
\item Exprimer, pour tout entier naturel $n$, $b_{n+1}$ en fonction de $b_n$.
\item En déduire la nature de la suite $(b_n)$ et préciser sa raison.
\end{enumerate}

\bigskip

\begin{minipage}{.60\linewidth}
\subsection*{PARTIE C : Comparaison des offres}

\bigskip

Pour tout entier naturel $n$ on note $c_n=a_n-b_n$.

\medskip

Les premiers termes des suites $(a_n)$, $(b_n)$ et $(c_n)$ sont donnés sur la feuille de calcul ci-contre.

\medskip

\begin{enumerate}
\item Quelle formule, destinée à être étirée vers le bas, doit-on écrire dans la cellule {\helvbx D2} ?
\item À partir de quelle année le montant de l'abonnement n°2 devient-il plus élevé que le montant de l'abonnement n°1 ?
\end{enumerate}
\end{minipage}\hfill
\begin{minipage}{.35\linewidth}\scriptsize\renewcommand{\arraystretch}{1.15}
\helvbx
\begin{tblr}{

	row{1} = {gray9,c},
	row{2} = {c},
	row{3-16} = {r},
	column{1} = {gray9,c},
	hlines,vlines,
	}
 & \hfill~A\hfill~ & \hfill~B\hfill~ & \hfill~C\hfill~ & \hfill~D\hfill~\\
1&$n$&\hfill~$a_n$\hfill~&\hfill~$b_n$\hfill~&\hfill~$c_n$\hfill~\\
2&0&250 \euro{}&200 \euro{}&50 \euro{}\\
3&1&280 \euro{}&220 \euro{}&60 \euro{}\\
4&2&310 \euro{}&242,00 \euro{}&68 \euro{}\\
5&3&340 \euro{}&266,20 \euro{}&73,80 \euro{}\\
6&4&370 \euro{}&292,82 \euro{}&77,18 \euro{}\\
7&5&400 \euro{}&322,10 \euro{}&77,90 \euro{}\\
8&6&430 \euro{}&354,31 \euro{}&75,69 \euro{}\\
9&7&460 \euro{}&389,74 \euro{}&70,26 \euro{}\\
10&8&490 \euro{}&428,72 \euro{}&61,28 \euro{}\\
11&9&520 \euro{}&471,59 \euro{}&48,41 \euro{}\\
12&10&550 \euro{}&518,75 \euro{}&31,25 \euro{}\\
13&11&580 \euro{}&570,62 \euro{}&9,38 \euro{}\\
14&12&610 \euro{}&627,69 \euro{}&$-17,69$ \euro{}\\
15&13&640 \euro{}&690,45 \euro{}&$-50,45$ \euro{}\\
16&14&670 \euro{}&759,50 \euro{}&$-89,50$ \euro{}\\
\end{tblr}
\end{minipage}

\end{document}
